\documentclass[12pt,a4paper,onecolumn]{article}

% Encoding, language, and layout per guide
\usepackage[utf8]{inputenc}
\usepackage[T1]{fontenc}
\usepackage[english]{babel}
\usepackage[margin=2.5cm]{geometry}
\usepackage{setspace}
\usepackage{newtxtext,newtxmath}
\usepackage{titlesec}
\usepackage{caption}
\usepackage{graphicx}
\usepackage{booktabs}
\usepackage{tabularx}
\usepackage{array}
\usepackage{ragged2e}
\usepackage{adjustbox}
\usepackage{float}
\usepackage{url}
\usepackage{hyperref}
\usepackage{enumitem}

\onehalfspacing
\setlength{\parindent}{0pt}
\setlength{\parskip}{6pt}

% Heading sizes per guide (12pt bold)
\titleformat{\section}{\bfseries\normalsize}{\thesection.}{0.6em}{}
\titleformat{\subsection}{\bfseries\normalsize}{\thesubsection}{0.6em}{}
\titlespacing*{\section}{0pt}{1.0\baselineskip}{0.6\baselineskip}
\titlespacing*{\subsection}{0pt}{0.8\baselineskip}{0.4\baselineskip}

% Captions per guide
\DeclareCaptionFont{tenpt}{\fontsize{10}{12}\selectfont}
\DeclareCaptionFont{ninept}{\fontsize{9}{11}\selectfont}
\captionsetup[table]{font={tenpt,bf},labelfont=bf,skip=6pt,justification=centering}
\captionsetup[figure]{font={ninept,bf},labelfont=bf,skip=6pt}
\setlength{\textfloatsep}{\baselineskip}
\setlength{\abovecaptionskip}{\baselineskip}
\setlength{\belowcaptionskip}{2\baselineskip}

% Metadata (replace with your actual info as needed)
\def\judul{MEGA-FrOG: Multi-Agent Enhanced Generation and Adaptation Framework of Open GraphRAG}
\def\judulInggris{MEGA-FrOG: Multi-Agent Enhanced Generation and Adaptation Framework of Open GraphRAG}
\def\penulisMahasiswa{Christopher Nathanael Wijaya; Febrian Dwi Kimhan; Muflih Naufal Maxi}
\def\penulisPembimbing{Dr.rer.nat. Fariz Darari, S.Kom., M.Sc.; Dr.techn. Fajar Juang Ekaputra, S.T., M.T.}
\def\afiliasiSatu{Department of Computer Science, Faculty of Computer Science, Universitas Indonesia}
\def\afiliasiDua{Department of Computer Science, Faculty of Computer Science, Universitas Indonesia}
\def\afiliasiTiga{Department of Computer Science, Faculty of Computer Science, Universitas Indonesia}
\def\afiliasiEmpat{Faculty of Computer Science, Universitas Indonesia}
\def\afiliasiLima{Faculty of Computer Science, Universitas Indonesia}
\def\emailSatu{christophernathanael@gmail.com, febriandwikimhan@gmail.com, muflihmaxi@gmail.com, fariz@ui.ac.id, fajar.ekaputra@wu.ac.at}

% Judul 14pt bold, centered (sesuai pedoman)
\makeatletter
\renewcommand\maketitle{%%
  \begin{center}
    {\fontsize{14}{17}\selectfont\bfseries \judul\par\vspace{12pt}}
    {\normalsize \penulisMahasiswa\par}
    {\normalsize \penulisPembimbing\par\vspace{12pt}}
    {\fontsize{10}{12}\selectfont
      1. \afiliasiSatu\\
      2. \afiliasiDua\\
      3. \afiliasiTiga\\
      4. \afiliasiEmpat\\
      5. \afiliasiLima\par}
    \vspace{12pt}
    {\fontsize{10}{12}\selectfont\itshape E-mail: \emailSatu\par}
  \end{center}
}
\makeatother

\begin{document}

\maketitle

% Abstract (EN)
\vspace*{2\baselineskip}
{\begin{center}\textbf{\fontsize{12}{14}\selectfont \judulInggris}\end{center}}
\vspace{12pt}
{\begin{center}\textbf{\fontsize{12}{14}\selectfont Abstract}\end{center}}
\vspace{12pt}
{\fontsize{10}{12}\selectfont\singlespacing
The rise of large language models (LLMs) has advanced information retrieval applications, but challenges such as knowledge updating difficulties, hallucinations, and lack of transparency persist. While Graph Retrieval-Augmented Generation (GraphRAG) addresses some issues, existing GraphRAG frameworks suffer from limited entity extraction capabilities, absence of specialized model training, and constrained Natural Language to SPARQL (NL2SPARQL) dataset generation. Our research presents MEGA-FrOG (Multi-Agent Enhanced Generation and Adaptation Framework of Open GraphRAG), an enhanced system addressing these limitations through three main improvements: comprehensive dataset generation using template-based and random-walk approaches, specialized model fine-tuning using QLoRA techniques on NL2SPARQL datasets, and sophisticated multi-agent architecture using LangGraph with intelligent routing and adaptive strategy selection. The enhanced pipeline integrates translation, entity extraction, strategy selection, verbalization, SPARQL generation, and answer generation agents with robust fallback mechanisms. The system is evaluated across four knowledge graphs: Wikidata, Curriculum KG, Legal KG, and GESIS KG, achieving Jaccard similarities of 0.318, 0.713, 0.487, and 0.546 respectively for the best-performing models. Comprehensive analysis reveals that specialized model fine-tuning provides the most significant performance improvements, with template-based dataset generation consistently outperforming random-walk approaches across all knowledge graph domains.\par}

{\fontsize{10}{12}\selectfont\singlespacing\textit{Keywords:} LLM, GraphRAG, Multi-Agent Systems, Knowledge Graphs, SPARQL, Model Fine-tuning \par}

\vspace*{2\baselineskip}
% {\begin{center}\textbf{\fontsize{12}{14}\selectfont \judulInggris}\end{center}}
\vspace{12pt}
{\begin{center}\textbf{\fontsize{12}{14}\selectfont Abstrak}\end{center}}
\vspace{12pt}
{\fontsize{10}{12}\selectfont\singlespacing
Munculnya \textit{large language models} (LLMs) telah memajukan aplikasi \textit{information retrieval}, tetapi tantangan seperti kesulitan pembaruan pengetahuan, halusinasi, dan kurangnya transparansi masih menjadi permasalahan. Meskipun \textit{Graph Retrieval-Augmented Generation} (GraphRAG) mengatasi beberapa masalah, kerangka kerja GraphRAG yang ada mengalami keterbatasan dalam kemampuan ekstraksi entitas, tidak adanya pelatihan (\textit{training}) model khusus, dan pembangkitan dataset \textit{Natural Language to SPARQL} (NL2SPARQL) yang terbatas. Penelitian ini memperkenalkan MEGA-FrOG (\textit{Multi-Agent Enhanced Generation and Adaptation Framework of Open GraphRAG}), sebuah sistem yang ditingkatkan untuk mengatasi keterbatasan tersebut melalui tiga perbaikan utama: pembuatan dataset komprehensif menggunakan pendekatan berbasis templat dan \textit{random-walk}, \textit{fine-tuning} model khusus menggunakan teknik QLoRA pada dataset NL2SPARQL, dan pembaruan arsitektur \textit{multi-agent} menggunakan LangGraph dengan penggunaan \textit{routing} dan pemilihan strategi adaptif. \textit{Pipeline} yang ditingkatkan mengintegrasikan translasi, ekstraksi entitas, pemilihan strategi, verbalisasi, pembuatan kueri SPARQL, dan pembangkitan jawaban dengan mekanisme \textit{fallback} yang lebih adaptif. Sistem dievaluasi pada empat graf pengetahuan: Wikidata, Curriculum KG, Legal KG, dan GESIS KG, mencapai similaritas Jaccard masing-masing 0,318, 0,713, 0,487, dan 0,546 untuk model dengan performa terbaik. Analisis komprehensif mengungkapkan bahwa \textit{fine-tuning} model khusus memberikan peningkatan performa yang paling signifikan, dengan pembangkitan dataset berbasis templat secara konsisten mengungguli pendekatan \textit{random-walk} di semua domain graf pengetahuan.\par}

{\fontsize{10}{12}\selectfont\singlespacing\textit{Kata kunci:} LLM, GraphRAG, Multi-Agent Systems, Knowledge Graphs, SPARQL, Model Fine-tuning \par}

\section*{Introduction}
Large language models (LLMs) have shifted IR applications toward prompt-driven QA and summarization. Despite high performance, LLMs face limitations in knowledge updating, interpretability, and hallucination. RAG augments parametric models with external knowledge so answers can be grounded in evidence, but chunk-based pipelines can lose context and lack explicit reasoning. Knowledge graphs encode semantic relations and enable graph reasoning. MEGA-FrOG leverages KGs to improve context quality, transparency, and accountability of QA. Our work extends the prior FrOG framework by introducing stronger graph-aware retrieval and verification while retaining practicality.

We address the following research questions:
\begin{enumerate}[leftmargin=*, label=RQ\arabic*:, itemsep=0.3\baselineskip]
  \item How to develop dataset generation methodologies to create diverse, high‑quality NL2SPARQL training data across multiple existing knowledge graph domains?
  \item How does specialized model fine‑tuning (QLoRA) improve the accuracy of entity/property extraction and SPARQL generation for direct KG querying compared to few‑shot prompting in the original GraphRAG framework?
  \item How can a multi‑agent (LangGraph‑based) architecture enhance adaptability, fault tolerance, and operational transparency compared to the original pipeline approach?
  \item How does the enhanced system perform in end‑to‑end QA across different existing knowledge graph domains and question complexity levels compared to the original baseline system?
\end{enumerate}

Our contributions are as follows:
\begin{enumerate}[leftmargin=*, label=C\arabic*), itemsep=0.3\baselineskip]
  \item Dual dataset generation methods (template‑based and random‑walk) for producing diverse, high‑quality NL2SPARQL datasets across Curriculum, Legal, and GESIS KGs, and enhancement of QALD‑9‑Plus for Wikidata with reasoning traces.
  \item Parameter‑efficient fine‑tuning (QLoRA) of specialized models for entity/property extraction and NL2SPARQL generation, replacing few‑shot prompting used in the original FrOG and yielding substantial performance gains.
  \item A LangGraph‑based multi‑agent architecture with runtime configurability, expanded fallbacks (including web search integration), and semantic logging to improve adaptability, fault tolerance, and transparency.
  \item Comprehensive end‑to‑end evaluation across four KGs and multiple complexity levels, quantifying improvements over FrOG and analyzing factors that affect performance (e.g., URI schemes, graph size; template vs. random‑walk training).
\end{enumerate}

\section*{Related Work}
\textbf{Retrieval-Augmented Generation.} RAG combines parametric models with non-parametric retrieval. Dense retrievers such as DPR learn representations for queries and passages; however, document chunking can degrade coherence and obscure cross-document relations. Variants introduce reranking, fusion-in-decoder, and hybrid BM25+dense pipelines to improve recall and precision, yet still operate over flat text.

\textbf{Knowledge Graphs for QA.} KGs model entities and relations as edges, enabling structured queries (SPARQL) and graph algorithms. For QA, KGs facilitate traceable evidence via subgraph walks and support disambiguation through typed relations. Recent works integrate KGs as either a retrieval source feeding LLMs or as a reasoning substrate with path-based justifications.

\textbf{GraphRAG and Graph-Aware Retrieval.} GraphRAG variants build or exploit a KG to provide higher-level summaries (communities, subgraphs) and relation-aware retrieval. Compared to text-only RAG, they better capture long-range dependencies and provide interpretable evidence chains.

\textbf{FrOG.} FrOG by Ongris and Tjitrahardja (2024) introduced an open GraphRAG pipeline driven largely by few-shot prompting and strong open-source integration. While practical, it faced limits in dataset diversity, reliance on prompting, and runtime configurability. MEGA-FrOG advances this line with dual dataset generation (template and random-walk), QLoRA fine-tuning for entity/property extraction and SPARQL generation, and a LangGraph-based multi-agent orchestration with richer fallbacks and semantic logging.

\section*{Materials and Methods}
\subsection*{System Overview}
The work enhances a KG‑centric QA system along three integrated components: (1) dataset generation (template‑based and random‑walk) to produce NL2SPARQL pairs with reasoning traces; (2) parameter‑efficient fine‑tuning (QLoRA) of models specialized for entity/property extraction and SPARQL generation; and (3) a LangGraph‑based multi‑agent architecture coordinating translation, extraction, strategy selection, SPARQL generation, and answer synthesis with configurable fallbacks and semantic logging. The system targets four existing knowledge graphs: Wikidata, Curriculum KG, Legal KG, and GESIS KG.

\subsection*{Dataset Generation}
We construct NL2SPARQL pairs using two complementary methods described in the thesis: (i) template-based generation that yields linguistically diverse questions aligned to schema patterns, and (ii) random-walk generation that discovers broader graph patterns from RDF data. Datasets include chain-of-thought style reasoning traces and are stratified by complexity (basic, intermediate, advanced), followed by SPARQL post-processing for formatting consistency.

\subsection*{Entity and Property Extraction}
Instead of rule-based extractors, we fine-tune LLMs (via QLoRA) to identify entities and properties mentioned in user questions. This specialized model provides more consistent extraction than few-shot prompting and aligns with the schemas of Wikidata, Curriculum, Legal, and GESIS KGs used in the experiments.

\subsection*{KG Construction and Storage}
Triples are materialized into a triple store capable of SPARQL queries and graph traversals. Indexing strategies (SPO/POS/OSP and quad variants) support efficient lookups. Named graphs encapsulate document-specific subgraphs to preserve provenance and enable subgraph-level operations.

\subsection*{SPARQL Generation and Execution}
A dedicated generation model (fine-tuned with QLoRA) translates natural language questions into SPARQL queries. Queries are executed against the KG (triple store) to retrieve answers and supporting triples. The system records intermediate steps for auditability and troubleshooting.

\subsection*{Agent Orchestration}
Using LangGraph, specialized agents coordinate translation, entity/property extraction, strategy selection, SPARQL generation, and answer synthesis with fallbacks (e.g., web search when configured). A controller governs routing and termination based on query complexity and user settings.

\subsection*{Design Decisions}
\begin{enumerate}[leftmargin=*, label=\arabic*), itemsep=0.4\baselineskip]
  \item \textbf{Dual dataset generation:} Combine template‑based questions for linguistic diversity with random‑walk generation for broader structural coverage.
  \item \textbf{Reasoning traces and stratification:} Include chain‑of‑thought style steps and stratify questions by complexity (basic, intermediate, advanced).
  \item \textbf{SPARQL post‑processing:} Normalize formatting for consistency with LLM outputs.
  \item \textbf{QLoRA fine‑tuning:} Train specialized models for entity/property extraction and NL2SPARQL generation instead of relying on few‑shot prompting.
  \item \textbf{LangGraph orchestration:} Employ configurable routing, fallbacks (e.g., web search), and semantic logging for transparency and robustness.
\end{enumerate}

\subsection*{Evaluation Protocol}
We conduct end‑to‑end evaluation using Jaccard Similarity on the generated and existing datasets across four knowledge graphs and multiple question complexity levels. We compare base instruction‑tuned models against their fine‑tuned counterparts (QLoRA) for both template‑based and random‑walk variants, and report average performance per domain. Experiments run on research‑grade hardware as described in the thesis setup.

\begin{figure}[htbp]
  \centering
  \includegraphics[width=0.9\linewidth]{assets/pdfs/pipeline2-cropped.pdf}
  \caption{MEGA-FrOG pipeline overview for extraction, KG storage, and retrieval--reasoning.}
\end{figure}

\section*{System Architecture}
The system is organized into three parts:
\begin{itemize}
  \item \textbf{Specialized models}: fine‑tuned (QLoRA) models for entity/property extraction and NL2SPARQL generation.
  \item \textbf{Orchestration}: a LangGraph workflow that calls the specialized models, builds SPARQL queries, executes them on a triple store (Jena/Fuseki), and collects evidence.
  \item \textbf{Web UI}: a simple interface to ask questions and view answers, SPARQL, and evidence.
\end{itemize}
The system works with four existing knowledge graphs: Wikidata, Curriculum, Legal, and GESIS.

\subsection*{Question Understanding and Answering}
The system extracts entities and properties from the user's question and then generates a SPARQL query. The triple store answers the query and returns rows and supporting triples. Reasoning is expressed directly in the SPARQL patterns (joins, filters, aggregates), which lets users inspect how the answer was derived.

\begin{figure}[htbp]
  \centering
  \includegraphics[width=0.85\linewidth]{assets/pics/kg_extractor_uml.jpg}
  \caption{UML-style overview of the KG extraction components and data flow.}
\end{figure}

\begin{figure}[htbp]
  \centering
  \includegraphics[width=0.75\linewidth]{assets/pics/langgraph_architecture_4.jpg}
  \caption{Reasoning and orchestration with graph-aware agents.}
\end{figure}

\subsection*{Backend and Services}
A backend exposes REST endpoints for asking questions and retrieving results. A LangGraph workflow coordinates the agents and records semantic logs. SPARQL queries are executed on Jena/Fuseki (TDB2). Web sockets stream progress and results to the UI.

\subsection*{User Interface}
The UI uses a chat-style layout. Users pick the target knowledge graph and model, ask a question, and then see: (i) the generated SPARQL, (ii) the answer, and (iii) the supporting triples. Simple controls allow switching graphs and toggling options.

\begin{figure}[htbp]
  \centering
  \includegraphics[width=0.9\linewidth]{assets/pics/homepage.png}
  \caption{System homepage with navigation, settings, and chat workspace.}
\end{figure}

\begin{figure}[htbp]
  \centering
  \includegraphics[width=0.95\linewidth]{assets/pics/jena_query_results.png}
  \caption{SPARQL query results interface for inspecting retrieved evidence.}
\end{figure}

\section*{Results}
\subsection*{Qualitative Findings}
The system works directly on existing knowledge graphs and returns a natural‑language answer grounded in SPARQL results, alongside the SPARQL and tabular output for verification. Across domains, fine‑tuned models produce queries that use the correct schema and direction of properties, while base or few‑shot approaches often misuse predicates or reverse subject/object roles.

Representative cases:
\begin{itemize}[leftmargin=*, itemsep=0.25\baselineskip]
  \item \textbf{Wikidata} (QALD‑9‑Plus): for “Who is the mayor of Berlin?”, the fine‑tuned model correctly uses \texttt{wd:Q64 wdt:P6 ?uri}, whereas the baseline variant sometimes selects an unrelated predicate.
  \item \textbf{Curriculum KG}: for “Alternative name for a course that uses Demo and has Programming Basics 1 as a prerequisite”, the fine‑tuned query selects the intended variable (\texttt{?alt\_name}) and predicates (e.g., \texttt{ns1:also\_known\_as}), while baseline variants produce noisier variable choices.
  \item \textbf{GESIS KG}: for “How many publications were published in 2000?”, the fine‑tuned query uses \texttt{schema:datePublished} correctly, while base models sometimes introduce non‑existent datatypes or properties.
  \item \textbf{Legal KG}: complex multi‑constraint questions remain challenging across systems; however, fine‑tuned models still adhere more closely to the domain schema than baselines.
\end{itemize}

\subsection*{Quantitative Findings}
Fine‑tuned models consistently outperform both base models and the FrOG baseline on Jaccard Similarity across all knowledge graphs. In Wikidata (QALD‑9‑Plus test set), LLaMA‑3.1‑8B jumps from near‑zero base (0.0079) to 0.3237, and Mistral‑Nemo improves from 0.0000 to 0.3024; Qwen2.5‑7B slightly surpasses FrOG (0.3176 vs. 0.3150). In Curriculum KG, template‑trained variants reach up to 0.7131, outperforming random‑walk training. In GESIS KG, gains are largest overall (average 0.4798; +33.3\% over FrOG). Legal KG also improves markedly (average 0.4216; +18.3\% over FrOG). Qwen2.5‑Coder‑7B attains the highest average after fine‑tuning. Notably, these gains are achieved without few‑shot prompts (FrOG relies on them).

% (Ablation-style analyses are not emphasized in the thesis; removed to reflect actual scope.)

\begin{figure}[htbp]
  \centering
  \includegraphics[width=0.8\linewidth]{assets/pics/graphrag_vs_rag.png}
  \caption{Schematic comparison of traditional RAG vs. GraphRAG/MEGA-FrOG information flow.}
\end{figure}

\begin{figure}[htbp]
  \centering
  \includegraphics[width=0.6\linewidth]{assets/pics/compar_new.jpg}
  \caption{Comparative performance summary across baseline and MEGA-FrOG conditions.}
\end{figure}

\subsection*{Results Summary}
Fine‑tuned models consistently outperform base models and the FrOG baseline across all knowledge graphs. Average Jaccard scores by domain: Wikidata 0.3101 (+4.9% over FrOG), Curriculum KG 0.6111 (+13.1%), Legal KG 0.4216 (+18.3%), and GESIS KG 0.4798 (+33.3%). Template‑trained variants generally surpass random‑walk training across domains. Qwen2.5‑Coder‑7B attains the highest average after fine‑tuning, and these gains are achieved without few‑shot prompts (whereas FrOG relies on them).

\begin{table}[H]
\centering
\renewcommand{\arraystretch}{1.15}
\setlength{\tabcolsep}{12pt}
\caption{Experiment Results on Wikidata (Jaccard Similarity)}
\label{table-wikidata-results}
\begin{adjustbox}{width=\textwidth}
  \begin{tabular}{| m{0.5\textwidth} | m{0.15\textwidth} | m{0.15\textwidth} | m{0.15\textwidth} |}
  \hline
  \RaggedRight{\textbf{Model}} & \textbf{FrOG} & \textbf{Base Model} & \textbf{Fine-tuned Model} \\
  \hline
  \RaggedRight{Qwen/Qwen2.5-7B-Instruct} & $0.3150$ & $0.2021$ & $\textbf{0.3176}$ \\
  \RaggedRight{Qwen/Qwen2.5-Coder-7B-Instruct} & $0.2907$ & $0.2531$ & $\textbf{0.2967}$ \\
  \RaggedRight{meta-llama/Llama-3.1-8B-Instruct} & $0.2884$ & $0.0079$ & $\textbf{0.3237}$ \\
  \RaggedRight{mistralai/Mistral-Nemo-Instruct-2407} & $0.2884$ & $0.0000$ & $\textbf{0.3024}$ \\
  \hline
  \end{tabular}
\end{adjustbox}
\vspace{1em}
\end{table}

\begin{table}[H]
\centering
\renewcommand{\arraystretch}{1.15}
\setlength{\tabcolsep}{12pt}
\caption{Experiment Results on Curriculum KG - Template Dataset (Jaccard Similarity)}
\label{table-curriculum-template-results}
\begin{adjustbox}{width=\textwidth}
  \begin{tabular}{| m{0.4\textwidth} | m{0.15\textwidth} | m{0.15\textwidth} | m{0.15\textwidth} | m{0.15\textwidth} |}
  \hline
  \RaggedRight{\textbf{Model}} & \textbf{FrOG} & \textbf{Base Model} & \textbf{Fine-tuned Template} & \textbf{Fine-tuned Random-Walk} \\
  \hline
  \RaggedRight{Qwen/Qwen2.5-7B-Instruct} & $0.5427$ & $0.5387$ & $\textbf{0.5963}$ & $0.4844$ \\
  \RaggedRight{Qwen/Qwen2.5-Coder-7B-Instruct} & $-$ & $0.4823$ & $\textbf{0.7042}$ & $0.5683$ \\
  \RaggedRight{meta-llama/Llama-3.1-8B-Instruct} & $-$ & $0.5211$ & $0.6423$ & $\textbf{0.7131}$ \\
  \RaggedRight{mistralai/Mistral-Nemo-Instruct-2407} & $-$ & $0.3516$ & $\textbf{0.6672}$ & $0.5217$ \\
  \hline
  \end{tabular}
\end{adjustbox}
\vspace{1em}
\end{table}

\begin{table}[H]
\centering
\renewcommand{\arraystretch}{1.15}
\setlength{\tabcolsep}{12pt}
\caption{Experiment Results on Curriculum KG - Random-Walk Dataset (Jaccard Similarity)}
\label{table-curriculum-random-results}
\begin{adjustbox}{width=\textwidth}
  \begin{tabular}{| m{0.4\textwidth} | m{0.15\textwidth} | m{0.15\textwidth} | m{0.15\textwidth} | m{0.15\textwidth} |}
  \hline
  \RaggedRight{\textbf{Model}} & \textbf{FrOG} & \textbf{Base Model} & \textbf{Fine-tuned Template} & \textbf{Fine-tuned Random-Walk} \\
  \hline
  \RaggedRight{Qwen/Qwen2.5-7B-Instruct} & $0.5382$ & $0.5133$ & $\textbf{0.6341}$ & $0.6127$ \\
  \RaggedRight{Qwen/Qwen2.5-Coder-7B-Instruct} & $-$ & $0.5333$ & $\textbf{0.6128}$ & $0.6047$ \\
  \RaggedRight{meta-llama/Llama-3.1-8B-Instruct} & $-$ & $0.5233$ & $\textbf{0.6033}$ & $0.5833$ \\
  \RaggedRight{mistralai/Mistral-Nemo-Instruct-2407} & $-$ & $0.4950$ & $\textbf{0.6251}$ & $0.6041$ \\
  \hline
  \end{tabular}
\end{adjustbox}
\vspace{1em}
\end{table}

\begin{table}[H]
\centering
\renewcommand{\arraystretch}{1.15}
\setlength{\tabcolsep}{12pt}
\caption{Experiment Results on GESIS KG - Template Dataset (Jaccard Similarity)}
\label{table-gesis-template-results}
\begin{adjustbox}{width=\textwidth}
  \begin{tabular}{| m{0.4\textwidth} | m{0.15\textwidth} | m{0.15\textwidth} | m{0.15\textwidth} | m{0.15\textwidth} |}
  \hline
  \RaggedRight{\textbf{Model}} & \textbf{FrOG} & \textbf{Base Model} & \textbf{Fine-tuned Template} & \textbf{Fine-tuned Random-Walk} \\
  \hline
  \RaggedRight{Qwen/Qwen2.5-7B-Instruct} & $0.3432$ & $0.2614$ & $\textbf{0.5328}$ & $0.4163$ \\
  \RaggedRight{Qwen/Qwen2.5-Coder-7B-Instruct} & $-$ & $0.2823$ & $\textbf{0.5463}$ & $0.4289$ \\
  \RaggedRight{meta-llama/Llama-3.1-8B-Instruct} & $-$ & $0.2742$ & $0.4526$ & $\textbf{0.5245}$ \\
  \RaggedRight{mistralai/Mistral-Nemo-Instruct-2407} & $-$ & $0.2537$ & $\textbf{0.5157}$ & $0.4083$ \\
  \hline
  \end{tabular}
\end{adjustbox}
\vspace{1em}
\end{table}

\begin{table}[H]
\centering
\renewcommand{\arraystretch}{1.15}
\setlength{\tabcolsep}{12pt}
\caption{Experiment Results on GESIS KG - Random-Walk Dataset (Jaccard Similarity)}
\label{table-gesis-random-results}
\begin{adjustbox}{width=\textwidth}
  \begin{tabular}{| m{0.4\textwidth} | m{0.15\textwidth} | m{0.15\textwidth} | m{0.15\textwidth} | m{0.15\textwidth} |}
  \hline
  \RaggedRight{\textbf{Model}} & \textbf{FrOG} & \textbf{Base Model} & \textbf{Fine-tuned Template} & \textbf{Fine-tuned Random-Walk} \\
  \hline
  \RaggedRight{Qwen/Qwen2.5-7B-Instruct} & $0.3764$ & $0.2403$ & $\textbf{0.4972}$ & $0.4872$ \\
  \RaggedRight{Qwen/Qwen2.5-Coder-7B-Instruct} & $-$ & $0.2654$ & $\textbf{0.4852}$ & $0.4776$ \\
  \RaggedRight{meta-llama/Llama-3.1-8B-Instruct} & $-$ & $0.2589$ & $0.4727$ & $\textbf{0.4927}$ \\
  \RaggedRight{mistralai/Mistral-Nemo-Instruct-2407} & $-$ & $0.2412$ & $\textbf{0.4793}$ & $0.4592$ \\
  \hline
  \end{tabular}
\end{adjustbox}
\vspace{1em}
\end{table}

\begin{table}[H]
\centering
\renewcommand{\arraystretch}{1.15}
\setlength{\tabcolsep}{12pt}
\caption{System Performance Comparison: Our Fine-tuned System vs. FrOG}
\label{table-system-comparison}
\begin{adjustbox}{width=\textwidth}
  \begin{tabular}{| m{0.25\textwidth} | m{0.15\textwidth} | m{0.15\textwidth} | m{0.15\textwidth} | m{0.15\textwidth} | m{0.15\textwidth} |}
  \hline
  \RaggedRight{\textbf{Knowledge Graph}} & \textbf{FrOG (Avg)} & \textbf{Our System (Fine-tuned Avg)} & \textbf{Absolute Improvement} & \textbf{Relative Improvement} & \textbf{Few-shot Usage} \\
  \hline
  \RaggedRight{Wikidata} & $0.2956$ & $\textbf{0.3101}$ & $+0.0145$ & $+4.9\%$ & FrOG only \\
  \RaggedRight{Curriculum KG} & $0.5404$ & $\textbf{0.6111}$ & $+0.0706$ & $+13.1\%$ & FrOG only \\
  \RaggedRight{Legal KG} & $0.3565$ & $\textbf{0.4216}$ & $+0.0651$ & $+18.3\%$ & FrOG only \\
  \RaggedRight{GESIS KG} & $0.3598$ & $\textbf{0.4798}$ & $+0.1200$ & $+33.3\%$ & FrOG only \\
  \hline
  \end{tabular}
\end{adjustbox}
\vspace{1em}
\end{table}

\subsection*{Additional Architectural and UI Figures}
\begin{figure}[H]
  \centering
  \includegraphics[width=0.85\linewidth]{assets/pics/react_comp.jpg}
  \caption{Componentized UI structure enabling maintainable front-end development.}
\end{figure}

\begin{figure}[H]
  \centering
  \includegraphics[width=0.95\linewidth]{assets/pics/febe_1.jpg}
  \caption{Frontend--backend interaction overview for evidence retrieval.}
\end{figure}

\begin{figure}[H]
  \centering
  \includegraphics[width=0.95\linewidth]{assets/pics/febe_2.jpg}
  \caption{End-to-end flow from query submission to answer rendering.}
\end{figure}

\section*{Discussion}
\subsection*{KG Characteristics and Performance}
Evaluation shows that knowledge graph characteristics strongly affect achievable performance. Domains with natural‑language URIs and more homogeneous schemas (e.g., Curriculum KG, GESIS KG) yield higher Jaccard scores than large, heterogeneous graphs with numeric identifiers (e.g., Wikidata).

\subsection*{Template vs. Random‑Walk Training}
Across all domains, template‑based fine‑tuning consistently outperforms random‑walk fine‑tuning, including when evaluated on diverse questions. This suggests that linguistically clear training data is more beneficial than broad structural coverage alone for NL2SPARQL.

\subsection*{System Integration Benefits}
Combining specialized models (extraction, SPARQL generation) with the multi‑agent orchestration improves end‑to‑end reliability compared to a few‑shot prompting baseline. Configurable fallbacks and semantic logging aid transparency and debugging.

\subsection*{Practical Notes}
The system runs on research‑grade hardware (e.g., Kaggle for training, local for evaluation) and is designed for reproducible deployment using open‑source stacks (triple store, web interface, orchestration).

\section*{Conclusion}
MEGA-FrOG demonstrates that KG-backed RAG improves QA accuracy and reliability while reducing hallucinations. We provide an architectural blueprint and empirical evidence supporting graph-aware retrieval and verification as practical enhancements to RAG.

\section*{Recommendations}
Future research should: (i) enhance extraction with supervised learning; (ii) optimize large-scale graph storage and indexing; (iii) extend evaluations to multiple domains and adversarial settings; and (iv) conduct user studies for usability and trust calibration.

\section*{References}
\begingroup\small
\begin{enumerate}
  \item Ongris, J. G., \& Tjitrahardja, E. (2024). FrOG: Framework of Open GraphRAG. Undergraduate Thesis, Universitas Indonesia, Faculty of Computer Science, Depok, Indonesia. Computer Science Study Program.
  \item Berners-Lee, T., Hendler, J., \& Lassila, O. (2001). The Semantic Web. \textit{Scientific American}. 
  \item Hogan, A., Blomqvist, E., Cochez, M., d'Amato, C., de Melo, G., Gutierrez, C., ... \& Zimmermann, A. (2021). Knowledge Graphs. \textit{arXiv:2003.02320}. \url{https://arxiv.org/abs/2003.02320}
  \item Wilcke, X., Bloem, P., \& De Boer, V. (2017). The knowledge graph as the default data model for learning on heterogeneous knowledge. \textit{Data Science, 1}(1-2), 39--57.
  \item Karpukhin, V., Oğuz, B., Min, S., Lewis, P., Wu, L., Edunov, S., Chen, D., \& Yih, W.-t. (2020). Dense Passage Retrieval for Open-Domain Question Answering. \textit{arXiv:2004.04906}. \url{https://arxiv.org/abs/2004.04906}
  \item Gao, Y., Xiong, Y., Gao, X., Jia, K., Pan, J., Bi, Y., ... \& Wang, H. (2024). Retrieval-Augmented Generation for Large Language Models: A Survey. \textit{arXiv:2312.10997}. \url{https://arxiv.org/abs/2312.10997}
  \item Peng, B., Zhu, Y., Liu, Y., Bo, X., Shi, H., Hong, C., Zhang, Y., \& Tang, S. (2024). Graph Retrieval-Augmented Generation: A Survey. \textit{arXiv:2408.08921}. \url{https://arxiv.org/abs/2408.08921}
  \item Zhang, Q., Chen, S., Bei, Y., Yuan, Z., Zhou, H., Hong, Z., ... \& Huang, X. (2025). A Survey of Graph Retrieval-Augmented Generation for Customized Large Language Models. \textit{arXiv:2501.13958}. \url{https://arxiv.org/abs/2501.13958}
  \item Vaswani, A., Shazeer, N., Parmar, N., Uszkoreit, J., Jones, L., Gomez, A. N., Kaiser, L., \& Polosukhin, I. (2017). Attention is All You Need. \textit{NeurIPS}. \url{https://proceedings.neurips.cc/paper_files/paper/2017/file/3f5ee243547dee91fbd053c1c4a845aa-Paper.pdf}
  \item Dettmers, T., Pagnoni, A., Holtzman, A., \& Zettlemoyer, L. (2023). QLoRA: Efficient Finetuning of Quantized LLMs. \textit{arXiv:2305.14314}. \url{https://arxiv.org/abs/2305.14314}
  \item Perevalov, A., Diefenbach, D., Usbeck, R., \& Both, A. (2022). QALD-9-Plus: A Multilingual Dataset for Question Answering over DBpedia and Wikidata. \textit{IEEE ICSC}, 229--234. \url{https://doi.org/10.1109/ICSC52841.2022.00045}
  \item Carroll, J. J., Dickinson, I., Dollin, C., Reynolds, D., Seaborne, A., \& Wilkinson, K. (2004). Jena: implementing the semantic web recommendations. In \textit{WWW Alt. Track Papers \& Posters} (pp. 74--83). ACM.
  \item Apache Software Foundation. (2023). Apache Jena Fuseki — A SPARQL Server. \url{https://jena.apache.org/documentation/fuseki2/}
  \item Apache Software Foundation. (2023). Apache Jena TDB2 — A native triple store for Jena. \url{https://jena.apache.org/documentation/tdb2/}
  \item Obitko, M. (2007). Semantic Web Architecture (Tutorial). \url{https://www.obitko.com/tutorials/ontologies-semantic-web/semantic-web-architecture.html}
  \item Raschka, S. (2024). \textit{Build a Large Language Model (From Scratch)}. Manning.
  \item OpenAI. (2023). GPT-4 Technical Report. \textit{arXiv:2303.08774}. \url{https://arxiv.org/abs/2303.08774}
\end{enumerate}
\endgroup

\end{document}
